% 05_numerical_experiments.tex
%
% Purpose:
%   Provide a standard paper-style placeholder for numerical experiments that
%   will be inserted after MATLAB validation.
%
% Main algorithm:
%   No numerical results are reported here.  The section lists planned
%   validation tests and experiment categories.
%
% Based on:
%   The validation and diagnostics discussion in the existing report draft.
%
% Main changes:
%   Adds a dedicated top-level experiments section without fabricating data,
%   figures, convergence tables, or numerical claims.
%
% Numerical goal:
%   Reserve space for reproducible experiments in the next revision.

\section{Numerical Experiments}

Numerical experiments will be inserted in the next revision after the MATLAB
drivers have been run and the outputs have been checked.  No numerical results
are reported in this draft.

\subsection{Validation tests and limiting cases}

Planned validation tests include an empty or homogeneous lead sanity check, a
perfect-crystal Bloch-mode recovery test, and limiting cases where the
line-defect formulation should reduce to a simpler periodic or empty-center
problem.

% TODO: Insert validated empty-cell and perfect-crystal comparison data.

\subsection{Convergence in boundary and Rayleigh discretizations}

The first convergence tables should vary the boundary discretization
\(N_{\mathrm{tot}}\), the Rayleigh truncation \(M\), proxy parameters, and
mode-selection thresholds.  The quantities to track include Bloch multipliers,
outgoing subspace projector differences, BIE residuals, port residuals, and
\(\sigmin(\Atr)\).

% TODO: Insert N_tot and M sweeps with residual and singular-value tables.

\subsection{Forced scattering examples}

Forced lead-scattering examples should prescribe incoming data from the selected
incoming trace subspaces and solve
\[
  \Atr x=b_{\mathrm{in}} .
\]
The planned outputs are field plots, reflected/transmitted outgoing
coefficients, and residual checks for the BIE and port equations.

% TODO: Insert forced scattering field plots and coefficient diagnostics.

\subsection{Eigenvalue and singular-value scans}

Eigenvalue or TEP-type examples should scan
\[
  \sigmin(\Atr(k,\beta))
\]
over selected ranges of \(k\) and \(\beta\), then refine candidate dips.  The
same candidates should be checked under changes in \(N_{\mathrm{tot}}\), \(M\),
and proxy parameters before being interpreted as physical modes.

% TODO: Insert singular-value scans and refined candidate tables.
